% English translation of chapter7.tex lines 1445--1660.
% Source: Methods of Algebra, Volume 2, commit
% 9a5803ff77dd3257484cb177f851a73770a59dd3 (CC BY 4.0).
% stable-unit-id: o014.aljabr2.chapter7.beck-monadicity-b
% Coverage: the remainder of Section 7.6, from the comparison functor through
% the linear version of Beck's theorem; Unit 096 begins Section 7.7.

\hypertarget{o014.aljabr2.chapter7.beck-monadicity-b}{ }
% segment-id: o014.aljabr2.chapter7.beck-monadicity-b.l001
\begin{lemma}\label{prop:T-comparison}
	Consider the monad $T$ and comonad $L$ arising from an adjoint pair $(F,G)$.
	Let $N\in\Obj(\mathcal{D})$. The data
	% segment-id: o014.aljabr2.chapter7.beck-monadicity-b.q001
	\[ M:=G(N)\in\Obj(\mathcal{C}), \quad
	 a:T(M)=GFG(N)\xrightarrow{G\varepsilon_N}G(N)=M \]
	give $(M,a)\in\Obj(\mathcal{C}^T)$. We thereby obtain a functor
	$\mathbb{K}:\mathcal{D}\to\mathcal{C}^T$ satisfying
	$G=U^T\mathbb{K}$.

	Dually, $F$ also has a canonical factorization
	$\mathcal{C}\xrightarrow{\mathbb{K}}\mathcal{D}^L\to\mathcal{D}$, which
	will henceforth be written $F=U^L\mathbb{K}$, where $U^L$ is the forgetful
	functor\footnote{Using the same symbol $\mathbb{K}$ will cause little
	confusion, since this book will not discuss monads and comonads
	simultaneously.}.
\end{lemma}
% segment-id: o014.aljabr2.chapter7.beck-monadicity-b.pf001
\begin{proof}
	It is enough to treat the case of $T$. In this case, the two diagrams in
	Definition~\ref{def:T-Alg} become
	% segment-id: o014.aljabr2.chapter7.beck-monadicity-b.g001
	\[\begin{tikzcd}
		GFGFG(N) \arrow[r, "GFG\varepsilon" inner sep=0.5em] \arrow[d, "{G\varepsilon FG}"'] & GFG(N) \arrow[d, "G\varepsilon"] \\
		GFG(N) \arrow[r, "G\varepsilon"'] & G(N)
	\end{tikzcd} \quad \begin{tikzcd}
		G(N) \arrow[r, "\eta G"] \arrow[equal, rd] & GFG(N) \arrow[d, "{G\varepsilon}"] \\
		& G(N).
	\end{tikzcd} \]
	Their commutativity is verified exactly as in
	Example~\ref{eg:adjunction-monad}, so it need not be repeated; the
	functoriality of $N\mapsto(M,a)$ is also immediate.
\end{proof}

% segment-id: o014.aljabr2.chapter7.beck-monadicity-b.p001
For a given adjoint pair $(F,G)$, applying a functor is a process that forgets
structure. Lemma~\ref{prop:T-comparison} factors $G$ (or $F$) as
$U^T\mathbb{K}$ (or $U^L\mathbb{K}$). We want to know whether structure is
lost only at the second step, when the forgetful functor $U^T$ is applied. If
so, the information lost in the process can be reconstructed with the aid of
the monad $T$ (or comonad $L$). This consideration motivates the following
concept.

% segment-id: o014.aljabr2.chapter7.beck-monadicity-b.d001
\begin{definition}\label{def:monadic}
	Consider a pair of adjoint functors
	% segment-id: o014.li2.chapter7.beck-monadicity-b.g002
	$\begin{tikzcd}
		F: \mathcal{C} \arrow[shift left, r] & \mathcal{D}: G \arrow[shift left, l]
	\end{tikzcd}$,
	which determines a monad $T$ on $\mathcal{C}$ and a comonad $L$ on
	$\mathcal{D}$.
	\begin{itemize}
		\item If the functor $\mathbb{K}:\mathcal{D}\to\mathcal{C}^T$ in
		Lemma~\ref{prop:T-comparison} is an equivalence of categories, the
		adjoint pair is called \emph{monadic}.
		\item If its dual version
		$\mathbb{K}:\mathcal{C}\to\mathcal{D}^L$ is an equivalence of
		categories, the adjoint pair is called \emph{comonadic}.
	\end{itemize}
\end{definition}

% segment-id: o014.aljabr2.chapter7.beck-monadicity-b.p002
For example, the free--forgetful adjoint pairs in
Examples~\ref{eg:Mon-monad} and \ref{eg:Mod-monad} are both monadic; a direct
check even shows that $\mathbb{K}$ is an isomorphism of categories in both
examples.

% segment-id: o014.aljabr2.chapter7.beck-monadicity-b.p003
Beck's theorem below, also called the monadicity theorem or the Barr--Beck
theorem, characterizes monadicity. We need some preparation.

% segment-id: o014.aljabr2.chapter7.beck-monadicity-b.dp001
\begin{definition-proposition}
	\index{split fork}
	Consider a diagram in a category $\mathcal{C}$
	% segment-id: o014.li2.chapter7.beck-monadicity-b.e001
	\begin{equation*}\label{eqn:split-coeq}\begin{tikzcd}[column sep=large]
		A \arrow[shift left, r, "u"] \arrow[shift right, r, "v" description] & B \arrow[r, "h"] \arrow[dashed, l, bend left, "t"] & Z \arrow[dashed, l, bend left, "s"]
	\end{tikzcd}\end{equation*}
	such that the solid part commutes, that is, $hu=hv$, and $h$ and $v$ have
	sections (that is, right inverses) $s$ and $t$, respectively, as indicated
	by the dashed arrows, with $sh=ut\in\End(B)$. The solid part of a diagram
	with these properties is called a \emph{split fork}. The diagram
	automatically gives a coequalizer of $u$ and $v$ in $\mathcal{C}$.
\end{definition-proposition}
% segment-id: o014.aljabr2.chapter7.beck-monadicity-b.pf002
\begin{proof}
	Consider the following diagram:
	% segment-id: o014.li2.chapter7.beck-monadicity-b.g003
	\[\begin{tikzcd}[row sep=small]
		& & W \\
		A \arrow[shift left, r, "u"] \arrow[shift right, r, "v"'] & B \arrow[r, "h"'] \arrow[ru, "k"] & Z
	\end{tikzcd} \qquad \begin{array}{ll}
		W \in \Obj(\mathcal{C}) \\
		ku = kv.
	\end{array} \]
	If there is a $\varphi:Z\to W$ with $\varphi h=k$, then
	$\varphi=\varphi hs=ks$. Conversely, $\varphi:=ks$ does satisfy
	$\varphi h=ksh=kut=kvt=k$. This proves the required universal property.
\end{proof}

% segment-id: o014.aljabr2.chapter7.beck-monadicity-b.p004
Unlike coequalizers in general, split forks are ``absolute'': their images
under arbitrary functors are still split forks.

% segment-id: o014.aljabr2.chapter7.beck-monadicity-b.eg001
\begin{example}\label{eg:T-split-fork}
	Let $(T,\mu,\eta)$ be a monad on $\mathcal{C}$. Every
	$(M,a)\in\Obj(\mathcal{C}^T)$ gives a split fork in $\mathcal{C}$:
	% segment-id: o014.li2.chapter7.beck-monadicity-b.e002
	\begin{equation*}\begin{tikzcd}[column sep=large]
			T^2(M) \arrow[shift left, r, "Ta"] \arrow[shift right, r, "\mu_M" description] & T(M) \arrow[r, "a"] \arrow[dashed, bend left, l, "\eta_{TM}"] & M \arrow[dashed, bend left, l, "\eta_M"]
		\end{tikzcd} \quad \begin{array}{ll}
			a (Ta) = a \mu_M , & \eta_M a = (Ta) \eta_{TM}, \\
			a \eta_M = \identity_M , & \mu_M \eta_{TM} = \identity_{TM}.
		\end{array}\end{equation*}
	The equalities on the right are part of Definitions~\ref{def:monad} and
	\ref{def:T-Alg}; for example,
	$\eta_Ma=(Ta)\eta_{TM}$ follows from the naturality of
	$\eta:\identity_{\mathcal{C}}\to T$. In particular, the diagram above
	realizes $M$ as the coequalizer of $Ta$ and $\mu_M$.
\end{example}

% segment-id: o014.aljabr2.chapter7.beck-monadicity-b.d002
\begin{definition}\label{def:Beck-split-pair}
	Consider a functor $G:\mathcal{D}\to\mathcal{C}$ and a pair of morphisms
	$f,g:X\rightrightarrows Y$ in $\mathcal{D}$. If there is a morphism
	$h:G(Y)\to Z$ such that the diagram
	% segment-id: o014.li2.chapter7.beck-monadicity-b.e003
	\begin{equation*}\begin{tikzcd}
			G(X) \arrow[shift left, r, "Gf"] \arrow[shift right, r, "Gg"'] & G(Y) \arrow[r, "h"] & Z
	\end{tikzcd}\end{equation*}
	is a split fork in $\mathcal{C}$, then $(f,g)$ is called a
	\emph{$G$-split pair} in $\mathcal{D}$.
\end{definition}

% segment-id: o014.aljabr2.chapter7.beck-monadicity-b.p005
Let $(f,g)$ be a $G$-split pair. Since $(Gf,Gg)$ has a coequalizer, it is
natural to ask whether that coequalizer can be lifted to a coequalizer of
$(f,g)$ in $\mathcal{D}$, and whether the lift is unique. A coequalizer is a
special case of a $\varinjlim$; the concepts of a functor creating or
preserving $\varinjlim$ in Definition~\ref{def:create-limit} provide the
right terminology for this question.

% segment-id: o014.aljabr2.chapter7.beck-monadicity-b.p006
Note that the forgetful functor $U^T:\mathcal{C}^T\to\mathcal{C}$ is
conservative in the sense of Definition~\ref{def:conservative}.

% segment-id: o014.aljabr2.chapter7.beck-monadicity-b.l002
\begin{lemma}\label{prop:UT-split-fork}
	Let $(T,\mu,\eta)$ be a monad on $\mathcal{C}$. The functor $U^T$ creates
	the coequalizer of every $U^T$-split pair.
\end{lemma}
% segment-id: o014.aljabr2.chapter7.beck-monadicity-b.pf003
\begin{proof}
	Suppose we are given a pair of morphisms in $\mathcal{C}^T$,
	$u,v:(M,a)\rightrightarrows(M',a')$, and a split fork in $\mathcal{C}$:
	% segment-id: o014.li2.chapter7.beck-monadicity-b.e004
	\begin{equation*}\begin{tikzcd}[column sep=large]
		M \arrow[shift left, r, "u"] \arrow[shift right, r, "v" description] & M' \arrow[r, "h"] \arrow[dashed, l, bend left, "t"] & M'' \arrow[dashed, l, bend left, "s"].
	\end{tikzcd}\end{equation*}
	Since this diagram gives a coequalizer in $\mathcal{C}$, it is enough to
	extend $M''$ uniquely to an object
	$(M'',a'')\in\Obj(\mathcal{C}^T)$ such that $h$ is a morphism in
	$\mathcal{C}^T$, and then to show that this construction gives the
	coequalizer of $u$ and $v$ in $\mathcal{C}^T$.

	The image of a split fork under any functor remains a split fork. Hence,
	in the following diagram, whose solid part is commutative,
	% segment-id: o014.li2.chapter7.beck-monadicity-b.g004
	\[\begin{tikzcd}
		T(M) \arrow[shift left, r, "Tu"] \arrow[shift right, r, "Tv"'] \arrow[d, "a"'] & T(M') \arrow[r, "Th"] \arrow[d, "{a'}"] & T(M'') \arrow[dashed, d, "{a''}"] \\
		M \arrow[shift left, r, "u"] \arrow[shift right, r, "v"'] & M' \arrow[r, "h"'] & M''
	\end{tikzcd}\]
	both rows are coequalizers in $\mathcal{C}$. The universal property
	therefore gives exactly one $a''$ that makes the whole diagram commute. If
	$(M'',a'')$ is shown to be a $T$-algebra, then commutativity of the
	right-hand part of this diagram says precisely that $h$ is a morphism in
	$\mathcal{C}^T$.

	For this, it remains to verify that the following two diagrams
	commute.\footnote{Translator's note: the lower-right node of the first
	diagram is printed as $T(M'')$ in the source. The unit law for a
	$T$-algebra requires that node to be $M''$; the target is corrected here.}
	% segment-id: o014.li2.chapter7.beck-monadicity-b.g005
	\[\begin{tikzcd}
		M'' \arrow[r, "{\eta_{M''}}"] \arrow[equal, rd] & T(M'') \arrow[d, "{a''}"] \\
		& M''
	\end{tikzcd} \quad
	\begin{tikzcd}
		T^2(M'') \arrow[r, "{\mu_{M''}}"] \arrow[d, "{Ta''}"'] & T(M'') \arrow[d, "{a''}"] \\
		T(M'') \arrow[r, "{a''}"'] & M'' .
	\end{tikzcd}\]
	The coequalizer diagrams show that $h$ and $Th$ are epimorphisms (as is
	$T^2h$). The commutativity of the two diagrams above therefore follows
	immediately from that of the corresponding diagrams for $M$ and $M'$.

	Finally, we show that $h$ gives the coequalizer of $u$ and $v$ in
	$\mathcal{C}^T$. Suppose there is a morphism in $\mathcal{C}^T$
	$k:(M',a')\to(W,b)$ with $ku=kv$.\footnote{Translator's note: the source
	prints the domain of $k$ as $(M'',a'')$, which makes $ku$, $kv$, and
	$\varphi h=k$ undefined. The correct domain is $(M',a')$.} In
	$\mathcal{C}$, there is a unique morphism $\varphi:M''\to W$ with
	$\varphi h=k$. It remains to prove that $\varphi$ is actually a morphism
	in $\mathcal{C}^T$. Consider the diagram
	% segment-id: o014.li2.chapter7.beck-monadicity-b.g006
	\[\begin{tikzcd}
		T(M') \arrow[r, "Th"] \arrow[d, "{a'}"'] & T(M'') \arrow[r, "T\varphi"] \arrow[d, "{a''}"] & T(W) \arrow[d, "b"] \\
		M' \arrow[r, "h"'] & M'' \arrow[r, "\varphi"'] & W.
	\end{tikzcd}\]
	The composites of the top and bottom rows are $Tk$ and $k$, respectively.
	Since $h$ and $k$ are morphisms in $\mathcal{C}^T$, the outer boundary and
	the left square commute. Thus the right square commutes after composition
	with $Th$. Since $Th$ is an epimorphism, the right square itself commutes.
	This shows that $\varphi$ is a morphism in $\mathcal{C}^T$.
\end{proof}

% segment-id: o014.aljabr2.chapter7.beck-monadicity-b.l003
\begin{lemma}\label{prop:Beck-FG-coeq}
	Let $G:\mathcal{D}\to\mathcal{C}$ be a functor with left adjoint $F$, and
	suppose $G$ creates the corresponding coequalizers for all $G$-split pairs
	$f,g:X\rightrightarrows Y$ (Definition~\ref{def:Beck-split-pair}). Then the
	following diagram is a coequalizer:
	% segment-id: o014.li2.chapter7.beck-monadicity-b.e005
	\begin{equation*}\begin{tikzcd}[column sep=large]
		FGFG(N) \arrow[shift left, r, "FG\varepsilon_N"] \arrow[shift right, r, "\varepsilon_{FG(N)}"'] & FG(N) \arrow[r, "\varepsilon_N"] & N,
	\end{tikzcd}\end{equation*}
	where $N\in\Obj(\mathcal{D})$.
\end{lemma}
% segment-id: o014.aljabr2.chapter7.beck-monadicity-b.pf004
\begin{proof}
	Applying $G$ to this diagram gives the split fork of
	Example~\ref{eg:T-split-fork} for
	$(M,a):=(G(N),G\varepsilon_N)=\mathbb{K}(N)$. Since $G$ creates the
	corresponding coequalizer, the original diagram is a coequalizer.
\end{proof}

% segment-id: o014.aljabr2.chapter7.beck-monadicity-b.t001
\begin{theorem}[J.\ M.\ Beck]\label{prop:Beck}
	\index{Beck's theorem}
	Let $G:\mathcal{D}\to\mathcal{C}$ be a functor with left adjoint $F$. The
	following statements are equivalent:
	\begin{enumerate}[(i)]
		\item the adjoint pair $(F,G)$ is monadic
		(Definition~\ref{def:monadic});
		\item $G$ creates the corresponding coequalizer for every $G$-split
		pair $f,g:X\rightrightarrows Y$;
		\item $G$ is conservative (Definition~\ref{def:conservative}), every
		$G$-split pair $f,g:X\rightrightarrows Y$ has a coequalizer
		$X\rightrightarrows Y\to Z$ in $\mathcal{D}$, and the image of that
		coequalizer under $G$ is a coequalizer of $(Gf,Gg)$.
	\end{enumerate}

	If any of the conditions above holds, a quasi-inverse functor $\mathbb{L}$
	to $\mathbb{K}:\mathcal{D}\to\mathcal{C}^T$ can be described as follows.
	For $(M,a)\in\Obj(\mathcal{C}^T)$, there is a coequalizer diagram
	% segment-id: o014.li2.chapter7.beck-monadicity-b.e006
	\begin{equation}\label{eqn:Beck-aux1}\begin{tikzcd}
		FGF(M) \arrow[shift left, r, "Fa"] \arrow[shift right, r, "\varepsilon_{FM}"'] & F(M) \arrow[r] & \mathbb{L}(M, a).
	\end{tikzcd}\end{equation}

	For the comonad determined by $(F,G)$, the criterion is entirely dual.
\end{theorem}
% segment-id: o014.aljabr2.chapter7.beck-monadicity-b.pf005
\begin{proof}
	First, we prove (i) $\implies$ (ii). Suppose
	$\mathbb{K}:\mathcal{D}\to\mathcal{C}^T$ has a quasi-inverse functor
	$\mathbb{L}$. The properties of being a split pair, being a split fork,
	and creating $\varinjlim$ are invariant under equivalences of categories.
	Since $G=U^T\mathbb{K}$, it is enough to prove that
	$U^T:\mathcal{C}^T\to\mathcal{C}$ creates the corresponding coequalizer
	for every $U^T$-split pair. This is precisely the content of
	Lemma~\ref{prop:UT-split-fork}.

	Next, (i) $\implies$ (iii). Since (i) $\implies$ (ii) is already known, it
	remains only to show that $G$ is conservative when $G$ is monadic. But
	$G=U^T\mathbb{K}$, where $\mathbb{K}$ is an equivalence of categories and
	$U^T$ is plainly conservative. Hence $G$ is conservative as well.

	We now prove (ii) $\implies$ (i) by constructing a quasi-inverse functor
	$\mathbb{L}$ to $\mathbb{K}$. For $(M,a)\in\Obj(\mathcal{C}^T)$,
	consider the pair of morphisms in $\mathcal{D}$
	% segment-id: o014.li2.chapter7.beck-monadicity-b.g007
	$\begin{tikzcd}
		FGF(M) \arrow[shift left, r, "Fa"] \arrow[shift right, r, "\varepsilon_{FM}"'] & F(M).
	\end{tikzcd}$
	Its image under $G$ is
	% segment-id: o014.li2.chapter7.beck-monadicity-b.g008
	$\begin{tikzcd}
		T^2(M) \arrow[shift left, r, "Ta"] \arrow[shift right, r, "\mu_M"'] & T(M)
	\end{tikzcd}$\!,
	which can be extended to a split fork by
	Example~\ref{eg:T-split-fork}. Thus $(Fa,\varepsilon_{FM})$ is a
	$G$-split pair, and the original diagram can be extended to a coequalizer
	diagram in $\mathcal{D}$, namely \eqref{eqn:Beck-aux1}.

	By \eqref{eqn:Beck-aux1} and the universal property of coequalizers, a
	morphism $(M,a)\xrightarrow{\varphi}(M',a')$ in $\mathcal{C}^T$ gives a
	unique morphism $\mathbb{L}(M,a)\to\mathbb{L}(M',a')$ making the following
	diagram commute:
	% segment-id: o014.li2.chapter7.beck-monadicity-b.e007
	\begin{equation*}\begin{tikzcd}
		FGF(M) \arrow[shift left, r, "Fa"] \arrow[shift right, r, "\varepsilon_{FM}"'] \arrow[d, "FGF\varphi"'] & F(M) \arrow[r] \arrow[d, "F\varphi"] & \mathbb{L}(M, a) \arrow[d] \\
		FGF(M') \arrow[shift left, r, "{Fa'}"] \arrow[shift right, r, "\varepsilon_{FM'}"'] & F(M') \arrow[r] & \mathbb{L}(M', a').
	\end{tikzcd}\end{equation*}
	Thus $(M,a)\mapsto\mathbb{L}(M,a)$ becomes a functor
	$\mathbb{L}:\mathcal{C}^T\to\mathcal{D}$.

	We next prove
	$\mathbb{L}\mathbb{K}\simeq\identity_{\mathcal{D}}$.\footnote{Translator's
	note: the source prints $\identity_{\mathcal{C}}$ here and in the
	comparison sentence below. Since $\mathbb{L}\mathbb{K}$ is an endofunctor
	of $\mathcal{D}$, the correct subscript is $\mathcal{D}$.} For every
	$N\in\Obj(\mathcal{D})$, the construction above gives a coequalizer
	diagram
	% segment-id: o014.li2.chapter7.beck-monadicity-b.e008
	\begin{equation*}\begin{tikzcd}[column sep=large]
		FGFG(N) \arrow[shift left, r, "FG\varepsilon_N"] \arrow[shift right, r, "\varepsilon_{FG(N)}"'] & FG(N) \arrow[r] & \mathbb{L}\mathbb{K}(N).
	\end{tikzcd}\end{equation*}

	Comparing this with the coequalizer diagram in
	Lemma~\ref{prop:Beck-FG-coeq} immediately gives the canonical isomorphism
	$\mathbb{L}\mathbb{K}\simeq\identity_{\mathcal{D}}$.

	Next, we verify
	$\mathbb{K}\mathbb{L}\simeq\identity_{\mathcal{C}^T}$. Clearly
	$\mathbb{K}F=\mathrm{Free}^T$; both send an object $M$ to
	$(GF(M),G\varepsilon_{FM})$. By definition and
	$\mu:=G\varepsilon F$, applying $\mathbb{K}$ to
	\eqref{eqn:Beck-aux1} gives
	% segment-id: o014.li2.chapter7.beck-monadicity-b.e009
	\begin{equation}\label{eqn:Beck-aux3}
		\begin{tikzcd}[row sep=small]
			(T^2(M), \ldots) \arrow[shift left, r, "Ta"] \arrow[shift right, r, "\mu_M"'] & (T(M), \ldots) \arrow[r] & \mathbb{K}\mathbb{L}(M, a). \\
			\mathrm{Free}^T(T(M)) \arrow[equal, u] & \mathrm{Free}^T(M) \arrow[equal, u] &
		\end{tikzcd}
	\end{equation}

	Applying the forgetful functor $U^T$ to the two terms on the left of
	\eqref{eqn:Beck-aux3} gives
	% segment-id: o014.li2.chapter7.beck-monadicity-b.g009
	$\begin{tikzcd}
		T^2(M) \arrow[shift left, r, "{Ta}"] \arrow[shift right, r, "{\mu_M}"'] & T(M)
	\end{tikzcd}$.
	Example~\ref{eg:T-split-fork} extends this diagram to a split fork with
	coequalizer $M$. On the other hand, $U^T$ creates this coequalizer
	(Lemma~\ref{prop:UT-split-fork}). Recalling
	Definition~\ref{def:create-limit}, it follows that
	\eqref{eqn:Beck-aux3} is also a coequalizer diagram.

	On the other hand, apply Lemma~\ref{prop:Beck-FG-coeq} to the adjoint pair
	$(\mathrm{Free}^T,U^T)$ and the object $N=(M,a)$. Carefully expanding the
	description of this adjoint pair gives the coequalizer diagram
	% segment-id: o014.li2.chapter7.beck-monadicity-b.g010
	\[\begin{tikzcd}[row sep=tiny]
		\mathrm{Free}^T(T(M)) \arrow[shift left, r, "Ta"] \arrow[shift right, r, "\mu_M"'] & \mathrm{Free}^T(M) \arrow[r] & (M, a).
	\end{tikzcd}\]
	Comparison with the preceding paragraph immediately gives
	$\mathbb{K}\mathbb{L}\simeq\identity_{\mathcal{C}^T}$.

	Finally, we prove (iii) $\implies$ (ii). The second part of (iii) says
	precisely that all $G$-split pairs have coequalizers and that $G$ preserves
	those coequalizers. Since $G$ is conservative,
	Remark~\ref{rem:conservative-reflect} shows that (ii) holds.

	Reversing all morphisms while retaining the direction of the functors,
	that is, replacing $\mathcal{C}$ and $\mathcal{D}$ by
	$\mathcal{C}^{\opp}$ and $\mathcal{D}^{\opp}$, gives the characterization
	of comonadicity. We omit the details.
\end{proof}

% segment-id: o014.aljabr2.chapter7.beck-monadicity-b.co001
\begin{corollary}
	Consider a pair of adjoint functors
	% segment-id: o014.li2.chapter7.beck-monadicity-b.g011
	$\begin{tikzcd}
		F: \mathcal{C} \arrow[shift left, r] & \mathcal{D}: G \arrow[shift left, l]
	\end{tikzcd}$
	with corresponding unit $\eta$ and counit $\varepsilon$. If this adjoint
	pair is monadic, then the diagram
	% segment-id: o014.li2.chapter7.beck-monadicity-b.g012
	\[\begin{tikzcd}[column sep=large]
		FGFG(N) \arrow[shift left, r, "FG\varepsilon_N"] \arrow[shift right, r, "\varepsilon_{FG(N)}"'] & FG(N) \arrow[r, "\varepsilon_N"] & N
	\end{tikzcd}\]
	is a coequalizer for every $N\in\Obj(\mathcal{D})$.
\end{corollary}
% segment-id: o014.aljabr2.chapter7.beck-monadicity-b.pf006
\begin{proof}
	Combine Lemma~\ref{prop:Beck-FG-coeq} with
	Theorem~\ref{prop:Beck}~(ii).
\end{proof}

% segment-id: o014.aljabr2.chapter7.beck-monadicity-b.p007
In practice, a linear version of Beck's theorem is often needed. More
precisely, let $\Bbbk$ be a commutative ring and $(F,G)$ an adjoint pair
between $\Bbbk\dcate{Mod}$-categories, with both functors $\Bbbk$-linear;
see Definition~\ref{def:cat-k-linear} and
Proposition~\ref{prop:automatic-k-morphism}. It is easy to see that
$\mathcal{C}^T$ naturally becomes a $\Bbbk\dcate{Mod}$-category and that,
by definition, both functors in
$\mathcal{D}\xrightarrow{\mathbb{K}}\mathcal{C}^T\xrightarrow{U^T}\mathcal{C}$
are $\Bbbk$-linear. If condition (ii) or (iii) of Beck's
Theorem~\ref{prop:Beck} is satisfied, the description by the diagram
\eqref{eqn:Beck-aux1} immediately shows that the quasi-inverse functor
$\mathbb{L}$ to $\mathbb{K}$ is also $\Bbbk$-linear. We therefore obtain a
$\Bbbk$-linear equivalence $\mathbb{K}:\mathcal{D}\to\mathcal{C}^T$.
